\scanpage{109}
La formule ci-dessus pour
\[
 \deleted{\unsure{d'où}}\alpha(n(x))=x\bar x
\]
donne, si on fait $x=y+z$
\[
\begin{aligned}
 \alpha\left(\underbrace{n(y+z)}_{\substack{\Vert\\n(y)+n(z)+\varphi_n(y,z)}}\right)
 &=\deleted{\unsure{\alpha(x\bar x)}}(y+z)\overline{(y+z)}=\\
 &=\underset{\substack{\Vert\\\alpha(n(y))}}{y\bar y}
 +\underset{\alpha(n(z))}{z\bar z}
 +\underbrace{y\bar z+\underbrace{z\bar y}_{=\overline{y\bar z}}}_{\tau(y\bar z)}
\end{aligned}
\]
donc
\[
 \alpha(\varphi_n(y,z))=\deleted{\unsure{\alpha t}}\tau(y\bar z)\qquad\text{i.e.}\quad\alpha(t(y\bar z))
\]
donc si $\alpha$ est injectif, on en déduit
\[
 \boxed{\varphi_n(y,z)=t(y\bar z)}
\]
i.e. la formule d'additivité pour $n$
\[
 \boxed{n(y+z)=n(y)+n(z)+\underset{\substack{\Vert\\t(\bar y z)}}{t(y\bar z)}}
 \qquad\text{car }t(x)=t(\bar x).
\]

Ceci posé, on veut

$3^\circ$) la formule précédente est vraie (sans supposer $\alpha$ injectif\ldots)

$4^\circ$) Une formule facultative, quand $k$ est fini loc. libre de rang \deleted{$\unsure{2}$} $n$ sur un anneau $k_0$, et \deleted{\unsure{alors on supposera}} $A$ loc. libre de rang $2n$ sur $k_0$ (N.B. donc $k_0\mathrel{-\!\!-\!\!-}A$ injectif, mais on ne suppose pas $k\mathrel{-\!\!-\!\!-}A$ injectif).
\[
 \boxed{\left\{\begin{aligned}
 \operatorname{Tr}_{k/k_0}(t(x))&=\operatorname{Tr}_{A/k_0}(x)\\
 N_{k/k_0}(n(x))&=N_{A/k_0}(x)
 \end{aligned}\right.}
\]
\margin{identité entre applic. \uline{linéaires} $A\mathrel{-\!\!-\!\!-}k_0$}
\margin{identité de $(2n)$-formes}
\mspage{108}
